% ============================================================
% Section 164: 有限深球方势阱的一般讨论
% ============================================================
\section{量子力学：有限深球方势阱的一般讨论}
\label{sec:finite-spherical-well}

\begin{modelbox}[题目：有限深球方势阱的一般讨论]
质量为 $\mu$ 的粒子在球方势阱
\[
V(r)=\begin{cases}0, & r<a \\ V_0, & r\ge a\end{cases}
\]
中运动，$V_0>0$。只考虑束缚态 ($0<E<V_0$)，设 $V_0$ 逐渐从小到大变化。
\begin{enumerate}
    \item 对每个角量子数 $l=0,1,2,\ldots$，求出现新束缚能级 ($E_{n,l}\approx V_0-0^+$) 的条件；
    \item 求出现第一个束缚态时 $V_0a^2$ 的值；
    \item 求各能级出现的先后次序；
    \item 若 $V_0$ 很大，估算束缚态总数 $N$。
\end{enumerate}
\end{modelbox}

\begin{tipbox}[考点快速分析]
\begin{itemize}
    \item \textbf{径向方程}：球 Bessel 函数 $j_l$ 和 $n_l$（阱内），衰减指数（阱外）
    \item \textbf{阱口条件}：$E\to V_0^-$ 时阱外方程化为 Euler 型，解为 $r^{-(l+1)}$
    \item \textbf{零点条件}：$j_{l-1}(k_0a)=0$ 决定新能级出现阈值
    \item \textbf{半经典估算}：相空间体积除以 $h^3$ 估算态总数
\end{itemize}
\end{tipbox}

\begin{derivationbox}[标准解题过程]
\textbf{1. 出现新束缚能级的条件}

束缚态波函数 $\psi_{nlm}=R_{nl}(r)Y_{lm}(\theta,\varphi)$。$E\approx V_0^-$ 时，$k_0=\sqrt{2\mu V_0}/\hbar$。

阱内（$r<a$）：$R''+\dfrac{2}{r}R'+\bigl[k_0^2-\dfrac{l(l+1)}{r^2}\bigr]R=0$，满足原点正则的解为 $R=j_l(k_0r)$。

阱外（$r\ge a$）：$E\approx V_0$ 使有效波数趋于零，方程化为 $R''+\dfrac{2}{r}R'-\dfrac{l(l+1)}{r^2}R=0$，衰减解为 $R\propto r^{-(l+1)}$。

$r=a$ 处 $R$ 与 $R'$ 连续。对 $R\propto r^{-(l+1)}$ 有 $\dfrac{d}{dr}[r^{l+1}R]=0$。因此阱内波函数在 $r=a$ 处须满足
\[
\frac{d}{dr}\bigl[r^{l+1}j_l(k_0r)\bigr]_{r=a}=0.
\]
利用递推关系 $\dfrac{d}{dx}[x^{l+1}j_l(x)]=x^{l+1}j_{l-1}(x)$，条件化为
\begin{equation}
\boxed{j_{l-1}(k_0a)=0,\quad l=0,1,2,\ldots}
\end{equation}
（$l=0$ 时 $j_{-1}(x)=\cos x/x$，条件为 $\cos k_0a=0$。）

\textbf{2. 第一个束缚态出现的条件}

由 Hellmann-Feynman 定理，给定 $l$ 下能量随 $V_0$ 增大而降低，且 $l$ 越小能级越深。第一个束缚态必为 $1s$ 态 ($l=0$)。

$l=0$ 时 $j_{-1}(k_0a)=\dfrac{\cos k_0a}{k_0a}=0$，第一根 $k_0a=\pi/2$。代入 $k_0$ 得
\begin{equation}
\boxed{V_0a^2=\frac{\pi^2\hbar^2}{8\mu}.}
\end{equation}

\textbf{3. 各能级出现的先后次序}

按 $j_{l-1}(x)$ 正零点的递增顺序，能级出现次序为
\begin{equation}
\boxed{1s,\;1p,\;1d,\;2s,\;1f,\;2p,\;1g,\;2d,\;3s,\;1h,\;2f,\;1i,\;3p,\;2g,\;1k,\;3d,\;4s,\;\ldots}
\end{equation}
（光谱符号：$s,p,d,f,g,h,i,k$ 对应 $l=0,1,2,3,4,5,6,7$。）

\textbf{4. 大 $V_0$ 时束缚态总数的估算}

半经典相空间估算：每个量子态占 $(2\pi\hbar)^3=h^3$。阱内坐标体积 $\dfrac{4\pi}{3}a^3$，动量上限 $p_0=\hbar k_0$，动量空间体积 $\dfrac{4\pi}{3}p_0^3$。

\begin{equation}
\boxed{N\approx\frac{2\pi^2}{9}\Bigl(\frac{ak_0}{\pi}\Bigr)^3=\frac{2\pi^2}{9}\Bigl(\frac{a}{\pi\hbar}\Bigr)^3(2\mu V_0)^{3/2}.}
\end{equation}
例如 $ak_0=7\pi/2$ 时，具体数出约 $99$ 个态（含简并），估算值约 $94$，吻合良好。
\end{derivationbox}

\begin{formulabox}[答案汇总]
\begin{center}
\begin{tabular}{ll}
\toprule
\textbf{结论} & \textbf{结果} \\
\midrule
新束缚态条件 & $j_{l-1}(k_0a)=0$（$k_0=\sqrt{2\mu V_0}/\hbar$） \\
第一个束缚态阈值 & $V_0a^2=\pi^2\hbar^2/(8\mu)$ \\
能级出现次序 & $1s,1p,1d,2s,1f,2p,1g,2d,3s,\ldots$ \\
大 $V_0$ 态总数 & $N\approx\dfrac{2\pi^2}{9}\left(\dfrac{ak_0}{\pi}\right)^3$ \\
\bottomrule
\end{tabular}
\end{center}
\end{formulabox}

\begin{insightbox}[物理图像]
\begin{itemize}
    \item 三维球方势阱与一维方势阱的关键区别：三维存在离心势垒 $l(l+1)\hbar^2/(2\mu r^2)$，使 $l$ 大的态更难束缚；一维对称势阱则无论多浅总至少有一个束缚态。
    \item $1s$ 态阈值 $V_0a^2=\pi^2\hbar^2/(8\mu)$ 恰与一维奇宇称态阈值相同，这是因为 $l=0$ 的径向方程在 $u(r)=rR(r)$ 变换后边界条件 $u(0)=0$ 与一维奇宇称 $x=0$ 处波函数为零类似。
    \item 深势阱极限下态数 $\propto(V_0a^2)^{3/2}$，这是三维相空间体积标度律的直接结果。
\end{itemize}
\end{insightbox}
